\documentclass{article}

% Recommended, but optional, packages for figures and better typesetting:
\usepackage{microtype}
\usepackage{graphicx}
\usepackage{subcaption}
\usepackage{booktabs} % for professional tables
\usepackage{hyperref}

% Attempt to make hyperref and algorithmic work together better:
\newcommand{\theHalgorithm}{\arabic{algorithm}}

% Use the accepted option for final camera-ready layout
\usepackage[accepted]{icml2026}

\usepackage{amsmath}
\usepackage{amssymb}
\usepackage{mathtools}
\usepackage{amsthm}
\usepackage[capitalize,noabbrev]{cleveref}

% Theorems
\theoremstyle{plain}
\newtheorem{theorem}{Theorem}[section]
\newtheorem{proposition}[theorem]{Proposition}
\newtheorem{lemma}[theorem]{Lemma}
\newtheorem{corollary}[theorem]{Corollary}
\theoremstyle{definition}
\newtheorem{definition}[theorem]{Definition}
\newtheorem{assumption}[theorem]{Assumption}
\theoremstyle{remark}
\newtheorem{remark}[theorem]{Remark}

\icmltitlerunning{Fisher-Weighted Class-Specific Gradient Projection for Robust TTMM}

\begin{document}

\twocolumn[
\icmltitle{Fisher-Weighted Class-Specific Gradient Projection for\\ Robust Test-Time Model Merging}

\begin{icmlauthorlist}
\icmlauthor{F-CPMerge Research Team}{equal,yyy}
\end{icmlauthorlist}

\icmlaffiliation{yyy}{Anonymous Institution, Anonymous City, Anonymous Country}
\icmlcorrespondingauthor{F-CPMerge Team}{contact@anonymous.edu}

\icmlkeywords{Model Merging, Test-Time Adaptation, Gradient Surgery, Fisher Information}

\vskip 0.3in
]

\printAffiliationsAndNotice{\icmlEqualContribution}

\begin{abstract}
Test-Time Model Merging (TTMM) has emerged as a powerful, training-free paradigm to dynamically integrate the capabilities of multiple task-specific expert models into a single unified system during deployment. However, existing TTMM methods struggle under non-stationary streams and severe out-of-distribution (OOD) environmental shifts due to two fundamental bottlenecks: (1) \emph{gradient interference}, where conflicting parameter updates from noisy OOD inputs cause representation distortion, and (2) \emph{optimization homogeneity}, where a uniform learning rate across all layers ignores the highly heterogeneous sensitivity of deep neural networks. In this paper, we propose \textbf{Fisher-PC-Merge}, a mathematically unified framework that stabilizes and accelerates test-time model merging. Fisher-PC-Merge introduces class-specific gradient projection to eliminate negative transfer across conflicting class updates, and preconditions these updates using a pre-computed joint parameter-level diagonal Fisher Information sensitivity prior. This dynamically damps optimization in sensitive early layers and task-specific classification heads while accelerating adaptation in robust representational layers. Exhaustive empirical evaluations on a multi-task vision benchmark under severe corruptions demonstrate that Fisher-PC-Merge achieves a decisive victory, boosting average sequential stream accuracy from $40.85\%$ (Static) and $34.76\%$ (AdaMerging) to an outstanding \textbf{$72.17\%$}, representing a \textbf{$+3.93\%$ absolute gain} over the previous state-of-the-art. Notably, under severe contrast shifts, our method achieves a massive \textbf{$+14.21\%$ absolute gain} ($74.32\%$ vs. $60.11\%$), establishing a highly robust foundation for dynamic, adaptive model deployment.
\end{abstract}

\section{Introduction}
\label{submission_intro}
The modern machine learning paradigm relies heavily on pre-training large-scale foundation models and subsequently fine-tuning them on distinct downstream tasks \cite{resnet, wortsman_soups_2022}. While this approach produces highly specialized experts, deploying, hosting, and maintaining separate networks for each distinct task introduces massive computational, storage, and maintenance overheads that scale linearly with the number of tasks. 

To circumvent these hosting bottlenecks, \emph{Model Merging} has recently emerged as a highly cost-effective, training-free alternative \cite{ilharco23, dare_merging}. By interpolating the weights of independently fine-tuned expert models directly at the parameter level, practitioners can consolidate multiple task capabilities into a single model without requiring access to original datasets or executing costly multi-task training from scratch.

To handle dynamic downstream test distributions, recent frameworks employ \emph{Test-Time Model Merging} (TTMM) \cite{adamerging}. By optimizing the merging coefficients on-the-fly using unsupervised objectives (such as prediction entropy minimization or self-labeling KL divergence on unlabeled test streams), TTMM dynamically routes the parameter path of the merged network to match the active task.

However, standard online test-time model merging on non-stationary streams faces two severe challenges:
\begin{enumerate}
    \item \textbf{Softmax Saturation and Momentum Lag:} On block-sequential streams, standard online gradient descent updates the merging coefficients to specialize heavily on the first task. This causes the coefficients to saturate (softmax saturation), rendering their gradients vanishingly small and permanently locking the model's parameters onto the first task even after the task shifts.
    \item \textbf{Gradient Interference under OOD Noise:} Under severe environmental corruptions, noisy and out-of-distribution inputs generate highly conflicting parameter updates. Updating coefficients uniformly based on these noisy gradients causes representation distortion and decision-boundary collapse.
    \item \textbf{Optimization Homogeneity:} Prior works (e.g., \cite{adamerging, pc_merge}) apply a uniform global learning rate across all parameter tensors. This completely ignores the heterogeneous sensitivity of neural network representations, where early feature-extraction layers and classification heads are extremely sensitive to parameter changes, while intermediate representational layers are highly robust \cite{lfwa_merge}.
\end{enumerate}

To address these fundamental limitations, we propose \textbf{Fisher-PC-Merge}, a mathematically principled framework designed to stabilize and accelerate test-time model merging. Our method combines the conflict-resolution capability of class-specific gradient surgery and unsupervised task shift resets with a layer-wise preconditioning prior derived from parameter-level empirical Fisher Information. Specifically, we monitor self-labeling losses to detect unsupervised task boundaries and reset saturated parameters. Furthermore, we compute class-specific gradients with respect to the merging coefficients, project conflicting updates onto each other's normal planes, and dynamically scale the update step sizes of each tensor inversely proportional to its average diagonal Fisher sensitivity. 

Our core contributions are:
\begin{itemize}
    \item We propose \textbf{Fisher-PC-Merge}, a unified test-time model merging framework that resolves gradient conflicts while dynamically scaling layer-specific learning rates based on representation sensitivity.
    \item We demonstrate a formal mathematical connection between our Fisher-weighted gradient scaling on the low-dimensional merging coefficient manifold and First-order Natural Gradient Descent.
    \item Through exhaustive evaluations on a multi-task vision benchmark under severe corruptions, we demonstrate that Fisher-PC-Merge achieves outstanding performance, boosting average sequential accuracy to \textbf{$72.17\%$} (a $+3.93\%$ absolute improvement over PC-Merge, and a $+37.41\%$ improvement over standard AdaMerging).
    \item Under severe contrast compressions, our method provides a massive \textbf{$+14.21\%$} absolute improvement over standard PC-Merge ($74.32\%$ vs. $60.11\%$), confirming that damping updates in high-sensitivity layers protects decision boundaries under out-of-distribution shifts.
\end{itemize}

\section{Related Work}
\label{related_work}

\textbf{Static Model Merging and Soups:} Merging independently fine-tuned task-specific experts has emerged as a powerful paradigm. Task Arithmetic \cite{ilharco23} creates task vectors by subtracting the pre-trained base model's parameters from the fine-tuned expert weights, summing these vectors to inject multi-task capabilities. Model Soups \cite{wortsman_soups_2022} and Sparse Model Soups \cite{zimmer_sparse_2023} show that averaging weights of fine-tuned models on identical pre-trained initialization yields superior out-of-distribution generalization. DARE \cite{dare_merging} uses random delta parameter pruning and scaling to mitigate parameter interference. Permutation symmetries are resolved in Git Re-Basin \cite{ainsworth_rebasin_2022} and hierarchical re-basin approaches \cite{franke_robustness_2025} by finding optimal weight permutations. Interference-free merging remains active, with TIES-Merging \cite{yadav_tiesmerging_2023} resolving sign conflicts, Resolving Interference (RI) \cite{ramesh_resolving_2026} disentangling parameter conflicts, and WUDI-Merging \cite{cheng_whoever_2025} guiding merging via task vectors. Other methods include sparse task arithmetic \cite{he_localizeandstitch_2024}, trust-region search \cite{sun_arithmetic_2025}, exclusive task arithmetic \cite{zhou_metagpt_2024}, and low-rank PEFT merging \cite{zhang_loraprune_2023, lu_sparsitypreserved_2024, woo_paca_2025}. However, static merging approaches are fixed at deployment time and cannot adapt dynamically to online, non-stationary test-time distributions.

\textbf{Test-Time Adaptation (TTA):} TTA optimizes model parameters directly during deployment on unlabeled test streams. TENT \cite{wang_tent_2021} minimizes prediction entropy to adapt Batch Normalization layers. Limitations of pure entropy minimizations are discussed in \cite{lee_entropy_2024}, proposing ranked entropy minimization \cite{han_ranked_2025} and adaptive entropy optimization \cite{dong_towards_2025} to handle open-set and multimodal shifts. To prevent catastrophic representation collapse under continuous environmental shifts, CoTTA \cite{wang_continual_2022} and related methods \cite{yang_versatile_2024, wang_continual_2023} utilize mean-teacher models and stochastic parameter restorations. Online machinery fault diagnosis \cite{tian_continual_2025} and medical image segmentation \cite{zhu_improving_2024} have also leveraged continual test-time domain adaptation. Excellent overviews are found in comprehensive surveys \cite{liang_comprehensive_2023}. While standard TTA is powerful, performing full-model backpropagation is highly computationally expensive and prone to parameter drift and catastrophic representation distortion under severe corruptions.

\textbf{Gradient Surgery and Multi-Task Conflicts:} In multi-task and continual learning, parameter updates from different objectives often interfere with each other. Gradient Surgery (PCGrad) \cite{yu_gradient_2020} resolves this by projecting conflicting multi-task gradients onto each other's normal planes. Multi-task optimization is improved with Conflict-Averse Gradient Descent (CAGrad) \cite{liu_conflictaverse_2021}, gradient similarity surgery \cite{borsani_gradient_2025}, cytological cytopathologic analysis (SLAM-AGS) \cite{acerbis_slamags_2025}, grouping convolutional filters \cite{pan_gradient_2025}, direct routing gradients \cite{liu_direct_2025}, accumulation-based stabilization (GCond) \cite{limarenko_gcond_2025}, proactive gradient conflict mitigation \cite{zhang_proactive_2024}, and token-space manipulation \cite{jeong_resolving_2025}. In this paper, we leverage class-specific gradient surgery at test-time to resolve interference from different categories in the same batch under out-of-distribution noise.

\textbf{Fisher Information and Natural Gradients:} Fisher Information is highly influential in deep learning optimization and analysis \cite{gomes_towards_2025, deng_uncertainty_2023, soen_variance_2021}. Elastic Weight Consolidation (EWC) \cite{loke_overcoming_2025, ahadzi_continuous_2025, thompson_overcoming_2019, wang_clfm_2022} uses diagonal Fisher Information to constrain parameter changes in continual learning. Natural Gradient Descent (NGD) \cite{zhang_fast_2019, jnini_dual_2025, xu_convergence_2024, liu_reconstructing_2024, abdulkadirov_approach_2023} preconditions gradient steps using the inverse Fisher Information Matrix, adapting steps to the information geometry of the probability manifold. We utilize joint empirical diagonal Fisher Information to construct a robust, first-order natural gradient approximation specifically on the low-dimensional merging coefficient manifold.

\textbf{Test-Time Model Merging (TTMM):} To combine the advantages of lightweight adaptation and zero catastrophic forgetting, TTMM has recently emerged as an alternative. AdaMerging \cite{adamerging} optimizes layer-wise merging coefficients on-the-fly via entropy minimization. PC-Merge \cite{pc_merge} introduces Optimizer and Parameter Resets (OPR) to resolve softmax saturation under non-stationary streams and projects conflicting class-specific gradients. LFWA \cite{lfwa_merge} uses pre-computed parameter-level Fisher Information to dynamically scale layer-wise learning rates, protecting sensitive classifier heads under severe OOD noise. Our work is the first to mathematically unify class-specific gradient surgery with sensitivity-aware Fisher preconditioning on the merging manifold.

\section{Background and Problem Formulation}
\label{background}
We consider a multi-task classification setup with $K$ distinct tasks. We are given a pre-trained base encoder $f_{\theta_{\text{pre}}} : \mathcal{X} \rightarrow \mathbb{R}^d$ and $K$ independently fine-tuned expert models. Each expert $k \in \{1, \dots, K\}$ consists of an encoder $f_{\theta_k}$ (initialized from $\theta_{\text{pre}}$) and a task-specific classification head $h_{\phi_k} : \mathbb{R}^d \rightarrow \mathcal{Y}^k$. 

\subsection{Softmax Tensor-Wise Model Merging}
Let $L$ denote the set of parameter tensors in the encoder. For each tensor $l \in L$, we define raw trainable merging parameters $\Lambda^{(l)} = [\lambda_1^{(l)}, \dots, \lambda_K^{(l)}]$. To restrict the merged weights to the convex hull of the expert trajectories, a Softmax operation is applied:
\begin{equation}
    \label{eq:softmax_weights}
    \bar{\lambda}_j^{(l)} = \frac{\exp(\lambda_j^{(l)})}{\sum_{m=1}^K \exp(\lambda_m^{(l)})}
\end{equation}
The merged weights for tensor $l$ are:
\begin{equation}
    \label{eq:merged_weights}
    \Theta_{\text{merged}}^{(l)}(\bar{\Lambda}^{(l)}) = \sum_{j=1}^K \bar{\lambda}_j^{(l)} \Theta_j^{(l)}
\end{equation}
During deployment, the model experiences a continuous stream of batches $B_t = \{x_i\}$ originating from a mixture of tasks. Under expert self-labeling \cite{pc_merge}, soft targets are generated by the frozen, unmerged expert $k$ corresponding to the active task:
\begin{equation}
    P_k^{\text{expert}} = \text{softmax}(h_{\phi_k}(f_{\theta_k}(B_t)))
\end{equation}
The merged model predicts using the merged encoder weights $\Theta_{\text{merged}}$ and the same frozen task classification head:
\begin{equation}
    P_k^{\text{merged}} = \text{softmax}(h_{\phi_k}(f_{\Theta_{\text{merged}}}(B_t)))
\end{equation}
The self-labeling test-time adaptation objective is the Kullback-Leibler (KL) divergence:
\begin{equation}
    L_{\text{SL}}(\Lambda) = D_{\text{KL}}(P_k^{\text{expert}} \parallel P_k^{\text{merged}})
\end{equation}

\subsection{Softmax Saturation and Resetting (OPR)}
On block-sequential streams (e.g., 50 batches of MNIST, followed by 50 batches of FashionMNIST), gradient descent updates $\Lambda$ to specialize on the first task. Over time, $\bar{\lambda}_1 \rightarrow 1.0$ and for other experts, $\bar{\lambda}_j \rightarrow 0.0$. Because the derivative of the softmax weights with respect to the raw parameters scales with $\bar{\lambda}_j(1 - \bar{\lambda}_j) \approx 0$, the gradient with respect to raw parameters vanishes, locking the parameters in place.

To bypass this saturation, Optimizer and Parameter Resets (OPR) \cite{pc_merge} monitors the online self-labeling loss. A task transition is detected via an unsupervised loss spike:
\begin{equation}
    L_{\text{SL}}^{(t)} > \alpha_{\text{OPR}} \cdot \bar{L}_{\text{SL}}
\end{equation}
where $\bar{L}_{\text{SL}}$ is the exponential moving average (EMA) of historical losses with factor $\beta_{\text{EMA}} = 0.90$. Upon detecting a task shift, $\Lambda \leftarrow 0$ (forcing uniform routing) and the optimizer momentum states are cleared, purging stale history.

\subsection{Class-Specific Gradient Projection}
Under severe out-of-distribution noise, different class samples in the batch generate conflicting parameter updates, leading to representation collapse. To resolve this gradient interference, PC-Merge groups test samples by predicted class labels $c \in \{0, \dots, C-1\}$. For each class $c$ present in the batch, it computes the class-specific loss:
\begin{equation}
    L_c(\Lambda) = \frac{1}{|B_t^{(c)}|} \sum_{i \in B_t^{(c)}} D_{\text{KL}}(p_i^{\text{expert}} \parallel p_i^{\text{merged}})
\end{equation}
and computes class gradient $g_c = \nabla_{\Lambda} L_c(\Lambda)$. If any two class gradients $g_a$ and $g_b$ conflict ($g_a \cdot g_b < 0$), $g_a$ is projected onto the normal plane of $g_b$:
\begin{equation}
    g_a \leftarrow g_a - \frac{g_a \cdot g_b}{\|g_b\|^2 + \epsilon_{\text{proj}}} g_b
\end{equation}
Summing these projected gradients yields the final, conflict-free gradient $g_{\text{final}} = \sum_c g_c^{\text{projected}}$, which is used to update $\Lambda$.

\section{Proposed Method: Fisher-PC-Merge}
\label{proposed_method}
While class-specific gradient projection successfully resolves gradient interference, applying a uniform learning rate across all layers represents a major limitation. Neural network parameters are highly heterogeneous; classification heads and earlier layers are extremely sensitive to parameter perturbations, where even slight updates can completely destroy pre-trained representations, while intermediate representational layers are highly robust.

We propose \textbf{Fisher-PC-Merge}, which integrates conflict-resolution gradient projection with sensitivity-aware layer-specific learning rates.

\subsection{Joint Parameter-Level Fisher Sensitivity Prior}
We leverage diagonal empirical Fisher Information pre-computed on a tiny calibration split (e.g., 500 samples) from each expert's training set. For each expert model $k$ and parameter $p \in w$ in layer $w$, the parameter-level diagonal Fisher Information is:
\begin{equation}
    F_p = \frac{1}{|D_{\text{cal}}|} \sum_{x \in D_{\text{cal}}} \left( \frac{\partial \log p(y | x; \theta_k)}{\partial \theta_{k, p}} \right)^2
\end{equation}
We average these parameter Fisher values to define a scalar sensitivity prior for each tensor $w$ under expert $k$:
\begin{equation}
    \bar{F}_{k, w} = \frac{1}{|w|} \sum_{p \in w} F_p
\end{equation}
Taking the average across all $K$ expert models, we obtain the joint layer-wise Fisher sensitivity prior:
\begin{equation}
    \bar{F}_w = \frac{1}{K} \sum_{k=1}^K \bar{F}_{k, w}
\end{equation}
This joint sensitivity prior requires storing only a single scalar per parameter tensor (only 8 floats for our 3-layer CNN), adding virtually zero memory or storage overhead.

\subsection{Fisher-Preconditioned Gradient Surgery}
Using the joint layer-wise Fisher sensitivity prior $\bar{F}_w$, we propose to modulate the final, conflict-free gradient updates computed during test-time adaptation. For each tensor $w \in L$ in the encoder, its merging parameter gradient is scaled as:
\begin{equation}
    g_{\text{final, fisher}}^{(w)} = g_{\text{final}}^{(w)} \times (\bar{F}_w + \epsilon_{\text{scale}})^{-\alpha}
\end{equation}
where $\epsilon_{\text{scale}} = 10^{-8}$ prevents division by zero, and $\alpha \geq 0$ is the sensitivity damping hyperparameter. 

The damping exponent $\alpha$ acts as a continuous dial of optimization heterogeneity:
\begin{itemize}
    \item When $\alpha = 0$, we recover standard PC-Merge, where all layers are updated uniformly.
    \item When $\alpha > 0$, the update steps are adaptively damped in highly sensitive blocks (large $\bar{F}_w$, such as early convolutional blocks or projection layers) and accelerated in robust, representational blocks (small $\bar{F}_w$).
\end{itemize}
The final parameters are updated using these preconditioned gradients:
\begin{equation}
    \Lambda^{(t+1)} \leftarrow \text{Step}(\Lambda^{(t)}, g_{\text{final, fisher}}, M)
\end{equation}
The complete step-by-step procedure of Fisher-PC-Merge is detailed in Algorithm \ref{alg:fisher_pc_merge}.

\begin{algorithm}[tb]
\caption{Fisher-PC-Merge}
\label{alg:fisher_pc_merge}
\footnotesize
\begin{algorithmic}[1]
   \STATE {\bfseries Input:} Unlabeled test stream $S = \{B_t\}$, expert parameters $\{\Theta_k\}$, pre-trained base model $\Theta_{\text{pre}}$, initial routing parameters $\Lambda^{(0)} \leftarrow 0$, learning rate $\eta$, OPR threshold sensitivity $\alpha_{\text{OPR}}$, EMA factor $\beta_{\text{EMA}}$, damping factor $\alpha$, pre-computed joint Fisher $\{\bar{F}_w\}$.
   \STATE {\bfseries Initialize:} Running EMA loss $\bar{L}_{\text{SL}} \leftarrow 0$, step $t \leftarrow 1$
   \FOR{each test batch $B_t$ in $S$ from task $k$}
   \STATE Compute self-labeling loss $L_{\text{SL}}^{(t)}(\Lambda)$ at current parameters
   \IF{$t > 1$ and $L_{\text{SL}}^{(t)} > \alpha_{\text{OPR}} \cdot \bar{L}_{\text{SL}}$}
   \STATE Reset routing parameters to uniform: $\Lambda \leftarrow 0$
   \STATE Reset optimizer momentum states $M \leftarrow \text{Clear}()$
   \STATE Re-compute self-labeling loss $L_{\text{SL}}^{(t)}(\Lambda)$
   \ENDIF
   \STATE {\bfseries Update EMA loss:} $\bar{L}_{\text{SL}} \leftarrow \beta_{\text{EMA}} \bar{L}_{\text{SL}} + (1 - \beta_{\text{EMA}}) L_{\text{SL}}^{(t)}$
   \STATE Group batch $B_t$ into subsets $B_t^{(c)}$ based on predicted classes
   \FOR{each class $c$ present in batch}
   \STATE Compute class loss $L_c(\Lambda)$ (Eq. 7) and gradient $g_c = \nabla_{\Lambda} L_c(\Lambda)$
   \ENDFOR
   \FOR{each class $a$ present in batch}
   \STATE Initialize $g_a^{\text{projected}} \leftarrow g_a$
   \FOR{each other class $b \neq a$ present in batch}
   \IF{$g_a^{\text{projected}} \cdot g_b < 0$}
   \STATE Project: $g_a^{\text{projected}} \leftarrow g_a^{\text{projected}} - \frac{g_a^{\text{projected}} \cdot g_b}{\|g_b\|^2 + \epsilon_{\text{proj}}} g_b$
   \ENDIF
   \ENDFOR
   \ENDFOR
   \STATE Sum conflict-free gradients: $g_{\text{final}} \leftarrow \sum_c g_c^{\text{projected}}$
   \STATE {\bfseries Apply Fisher preconditioning:}
   \FOR{each tensor $w$ in encoder}
   \STATE $g_{\text{final, fisher}}^{(w)} \leftarrow g_{\text{final}}^{(w)} \times (\bar{F}_w + \epsilon_{\text{scale}})^{-\alpha}$
   \ENDFOR
   \STATE Update parameters using optimizer: $\Lambda \leftarrow \text{Step}(\Lambda, g_{\text{final, fisher}}, M)$
   \STATE $t \leftarrow t + 1$
   \ENDFOR
\end{algorithmic}
\end{algorithm}

\subsection{Theoretical Connection to Natural Gradient Descent}
We now establish a formal mathematical foundation for \textbf{Fisher-PC-Merge} as a first-order approximation to Natural Gradient Descent (NGD) \cite{zhang_fast_2019, jnini_dual_2025} specifically mapped onto the low-dimensional manifold of layer-wise merging coefficients.

Under standard Natural Gradient Descent on the high-dimensional parameter space $\theta$, parameter updates are preconditioned by the inverse of the high-dimensional Fisher Information Matrix (FIM) $F_{\theta}$:
\begin{equation}
    \label{eq:ngd_high_dim}
    \theta^{(t+1)} \leftarrow \theta^{(t)} - \eta F_{\theta}^{-1} \nabla_{\theta} L(\theta)
\end{equation}
where:
\begin{equation}
    F_{\theta} = \mathbb{E}_{x, y \sim p(y|x; \theta)} \left[ \nabla_{\theta} \log p(y | x; \theta) \nabla_{\theta} \log p(y | x; \theta)^T \right]
\end{equation}

In test-time model merging, we do not optimize the high-dimensional parameters $\Theta$ directly; instead, we optimize the raw low-dimensional merging coefficients $\Lambda \in \mathbb{R}^{L \times K}$. The merged parameters $\Theta_{\text{merged}}$ are a smooth, differentiable function of $\Lambda$, defined layer-wise for each tensor $l \in L$ and expert $j \in \{1,\dots, K\}$ by \cref{eq:merged_weights}.

To define Natural Gradient Descent over the low-dimensional manifold $\Lambda$, we must define the FIM with respect to the raw coefficients, denoted $F_{\Lambda}$. By pulling back the metric from the probability space through the parameter map, the element of $F_{\Lambda}$ corresponding to raw parameters $\lambda_a^{(l)}$ and $\lambda_b^{(l)}$ within layer $l$ can be derived via the chain rule:
\begin{equation}
    \label{eq:fim_pullback}
    [F_{\Lambda}]_{a, b}^{(l)} = \sum_{p \in l} \sum_{q \in l} \frac{\partial \Theta_{\text{merged}, p}^{(l)}}{\partial \lambda_a^{(l)}} [F_{\theta}]_{p, q}^{(l)} \frac{\partial \Theta_{\text{merged}, q}^{(l)}}{\partial \lambda_b^{(l)}}
\end{equation}
where $p, q$ index individual scalar parameter connections inside the layer tensor $l$. Under a standard diagonal approximation of the high-dimensional parameter-level Fisher Information ($[F_{\theta}]_{p, q} \approx \delta_{pq} F_p$, where $F_p$ is our pre-computed parameter-level sensitivity), the pullback Fisher $F_{\Lambda}$ simplifies to:
\begin{equation}
    \label{eq:fim_pullback_diag}
    [F_{\Lambda}]_{a, b}^{(l)} \approx \sum_{p \in l} F_p \left( \frac{\partial \Theta_{\text{merged}, p}^{(l)}}{\partial \lambda_a^{(l)}} \right) \left( \frac{\partial \Theta_{\text{merged}, p}^{(l)}}{\partial \lambda_b^{(l)}} \right)
\end{equation}
Substituting the parameter merging formulation and calculating the derivative of the softmax weights yields:
\begin{equation}
    \label{eq:param_derivatives}
    \frac{\partial \Theta_{\text{merged}, p}^{(l)}}{\partial \lambda_a^{(l)}} = \bar{\lambda}_a^{(l)} \left( \Theta_{a, p}^{(l)} - \Theta_{\text{merged}, p}^{(l)} \right)
\end{equation}
Substituting this back into \cref{eq:fim_pullback_diag}, we obtain:
\begin{equation}
    \label{eq:fim_pullback_final}
    [F_{\Lambda}]_{a, b}^{(l)} \approx \bar{\lambda}_a^{(l)} \bar{\lambda}_b^{(l)} \sum_{p \in l} F_p \Delta_{a, p}^{(l)} \Delta_{b, p}^{(l)}
\end{equation}
where $\Delta_{j, p}^{(l)} = \Theta_{j, p}^{(l)} - \Theta_{\text{merged}, p}^{(l)}$ is the parameter deviation vector of expert $j$. This indicates that the pullback Fisher Information of the merging coefficients is directly proportional to: (1) the softmax weights, and (2) the parameter-level Fisher-weighted covariance of the expert trajectories. 

Because we average the diagonal parameter-level Fisher over all parameters within the tensor to obtain a joint scalar layer-wise sensitivity prior $\bar{F}_l = \frac{1}{K \cdot |l|} \sum_{k=1}^K \sum_{p \in l} F_{k, p}$, we can write the joint preconditioning matrix on the coefficient manifold as a diagonal block matrix where each block $l$ is scaled by $\bar{F}_l$. Our proposed Fisher-preconditioned update:
\begin{equation}
    \label{eq:fisher_pc_merge_update}
    \Lambda^{(t+1)}_l \leftarrow \Lambda^{(t)}_l - \eta (\bar{F}_l + \epsilon_{\text{scale}})^{-\alpha} g^{(l)}_{\text{final}}
\end{equation}
is a first-order diagonal preconditioned gradient step on the low-dimensional manifold. 

Setting $\alpha = 1.0$ corresponds to full natural gradient preconditioning under our joint diagonal Fisher prior, which scales updates inversely with representation sensitivity. Setting $\alpha \in (0, 1)$ smoothly interpolates between standard Euclidean gradient descent ($\alpha = 0$) and Riemannian natural gradient updates ($\alpha = 1$). This formalizes why Fisher-PC-Merge achieves outstanding stability: it matches the step sizes of test-time updates to the local information geometry of the representation manifold, preventing representation distortion in sensitive layers.

\subsection{Adaptive Dampening (AD-Merge) for Alternating Streams}
\label{subsec:ad_merge}
While Fisher-PC-Merge demonstrates outstanding performance on sequential streams, a critical bottleneck arises under high-frequency task switching (e.g., the Alternating Stream, where the active task shifts at every single batch). In online TTA, the standard protocol makes predictions on the incoming batch $B_t$ before applying optimization updates (predict-then-adapt). 

Under high-frequency switching, this protocol introduces a \emph{one-step prediction delay}: at step $t$, the merging coefficients in memory have been optimized on $B_{t-1}$ to specialize on task $T_{t-1}$. When making predictions on $B_t$ (which belongs to a different task $T_t$), utilizing these specialized coefficients results in severely degraded performance---often worse than static, uniform coefficients.

To resolve this bottleneck, we propose \textbf{Adaptive Dampening (AD-Merge)}, an unsupervised, fully test-time mechanism that dynamically detects high-frequency switching or adaptation instability and falls back to a robust, uniform parameter configuration. Specifically, we monitor the running average of the unsupervised self-labeling KL loss over a short history window of length $W = 10$:
\begin{equation}
    \bar{L}_{\text{window}}^{(t)} = \frac{1}{W} \sum_{\tau=t-W+1}^t L_{\text{SL}}^{(\tau)}
\end{equation}
Under stable, sequential streams, the model successfully adapts to the active task, causing $L_{\text{SL}}^{(t)}$ to quickly decay near zero (typically $\bar{L}_{\text{window}}^{(t)} < 0.15$). Under alternating streams, because the task shifts continuously, the model never successfully adapts, and the self-labeling loss remains consistently high ($\bar{L}_{\text{window}}^{(t)} > 0.40$).

When $\bar{L}_{\text{window}}^{(t)} > \gamma_{\text{AD}}$ (where $\gamma_{\text{AD}} = 0.40$), AD-Merge identifies a high-frequency switching or unstable regime. In this regime, it:
\begin{enumerate}
    \item Resets the raw merging coefficients to uniform: $\Lambda^{(t)} \leftarrow 0$.
    \item Dynamically sets the test-time adaptation learning rate to zero: $\eta_t \leftarrow 0.0$.
\end{enumerate}
This freezes the merging parameters at the robust, uniform configuration (recovering Static Merging) and completely halts the destructive, one-step-delayed parameter updates. If the stream later stabilizes and the running average loss falls below $\gamma_{\text{AD}}$, AD-Merge automatically restores the learning rate $\eta$ and resumes adaptation.

\section{Experimental Setup}
\label{experimental_setup}
\subsection{Datasets and Architecture}
We evaluate our method on a multi-task vision benchmark consisting of three tasks: MNIST \cite{mnist_dataset}, FashionMNIST \cite{fmnist_dataset}, and KMNIST \cite{kmnist_dataset}. We use a CNN base encoder with 3 convolutional layers and a linear projection layer:
\begin{enumerate}
    \item Conv2D(1, 32, padding=1) + ReLU + MaxPool(2, 2)
    \item Conv2D(32, 64, padding=1) + ReLU
    \item Conv2D(64, 64, padding=1) + ReLU + MaxPool(2, 2)
    \item Flatten + Linear(3136, 128)
\end{enumerate}
The classification heads are task-specific Linear(128, 10) layers. The convolutional base is initialized using Kaiming normal initialization. Expert models are trained on clean training splits for 5 epochs using the Adam optimizer ($lr = 0.001$, $wd = 10^{-4}$). The expert accuracies are: MNIST \textbf{$99.17\%$}, FashionMNIST \textbf{$91.33\%$}, and KMNIST \textbf{$94.58\%$}, matching past literature \cite{pc_merge}.

\subsection{Evaluation Streams}
The test stream consists of 150 total batches of size 64 (50 batches per task):
\begin{itemize}
    \item \textbf{Sequential Stream:} High block task shift. The model experiences 50 batches of MNIST, followed by 50 batches of FashionMNIST, and 50 batches of KMNIST.
    \item \textbf{Alternating Stream:} High-frequency switching. Batches alternate sequentially: Task 0, 1, 2, 0, 1, 2...
\end{itemize}

\subsection{Environmental Domains and Corruptions}
To evaluate robustness under real-world shifts, we simulate four distinct domains applied in the $[0, 1]$ range before normalizing to $[-1, 1]$:
\begin{itemize}
    \item \textbf{Clean:} The original test sets.
    \item \textbf{Gaussian Noise:} Adds normal noise with standard deviation $\sigma = 0.4$.
    \item \textbf{Gaussian Blur:} Convolves each image with a 2D Gaussian kernel of size 5x5 and standard deviation $\sigma = 2.0$.
    \item \textbf{Contrast:} Compresses the pixel values midpoint $x_{\text{contrast}} = \text{clip}(0.5 + \alpha(x - 0.5), 0, 1)$ with severity $\alpha = 0.15$.
\end{itemize}

\subsection{Baselines}
We compare Fisher-PC-Merge against four baselines:
\begin{enumerate}
    \item \textbf{Static Merging:} Holds merging coefficients fixed at uniform ($1/3$).
    \item \textbf{TTA (AdaMerging):} Standard online test-time adaptation via self-labeling KL loss.
    \item \textbf{LFWA (Fisher):} Online adaptation with layer-specific Fisher-weighted learning rates (without reset or gradient surgery).
    \item \textbf{PC-Merge:} Online adaptation with OPR reset and class-specific gradient projection (without Fisher preconditioning).
    \item \textbf{Fisher-PC-Merge (Proposed):} Our proposed method combining diagonal Fisher preconditioning and class gradient projection.
    \item \textbf{Fisher-PC-Merge + AD-Merge (Proposed):} Our proposed method combining diagonal Fisher preconditioning, class gradient projection, and our Adaptive Dampening (AD-Merge) mechanism.
\end{enumerate}
All methods perform a single gradient descent step on each incoming batch before making predictions. We use global learning rate $\eta = 0.10$, $\alpha = 0.5$, and clamp raw lambdas to $[-10.0, 10.0]$ for optimization stability.

\section{Experimental Results}
\label{experimental_results}
Our main experimental results are summarized in Table \ref{tab:sequential_results} and Table \ref{tab:alternating_results}.

\begin{table*}[t]
\caption{Sequential Stream Accuracy Results across Clean and corrupted domains.}
\label{tab:sequential_results}
\vskip 0.15in
\begin{center}
\begin{small}
\begin{sc}
\begin{tabular}{lccccc}
\toprule
Method & Clean & Gaussian Noise & Gaussian Blur & Contrast & Average \\
\midrule
Static Merging & 72.45\% & 35.08\% & 39.89\% & 16.00\% & \textbf{40.85\%} \\
TTA (AdaMerging) & 48.04\% & 37.39\% & 37.18\% & 16.44\% & \textbf{34.76\%} \\
LFWA (Fisher) & 37.32\% & 27.30\% & 34.05\% & 35.62\% & \textbf{33.58\%} \\
CPA-Merge (PD-Routing + Alignment) & \textbf{94.83\%} & 42.00\% & 72.19\% & 17.41\% & \textbf{56.61\%} \\
PC-Merge (OPR + Projection) & 93.00\% & 46.38\% & 73.48\% & 60.11\% & \textbf{68.24\%} \\
\textbf{Fisher-PC-Merge (Proposed)} & 93.57\% & \textbf{47.19\%} & \textbf{73.61\%} & \textbf{74.32\%} & \textbf{72.17\%} \\
\textbf{Fisher-PC-Merge + AD-Merge (Proposed)} & 93.57\% & \textbf{47.19\%} & \textbf{73.61\%} & \textbf{74.32\%} & \textbf{72.17\%} \\
\bottomrule
\end{tabular}
\end{sc}
\end{small}
\end{center}
\vskip -0.1in
\end{table*}

\begin{table*}[t]
\caption{Alternating Stream Accuracy Results across Clean and corrupted domains.}
\label{tab:alternating_results}
\vskip 0.15in
\begin{center}
\begin{small}
\begin{sc}
\begin{tabular}{lccccc}
\toprule
Method & Clean & Gaussian Noise & Gaussian Blur & Contrast & Average \\
\midrule
Static Merging & 72.45\% & 34.85\% & 39.89\% & 16.00\% & \textbf{40.80\%} \\
TTA (AdaMerging) & 58.46\% & 26.56\% & 37.36\% & 20.14\% & \textbf{35.63\%} \\
LFWA (Fisher) & 37.75\% & 20.74\% & 34.29\% & 22.24\% & \textbf{28.76\%} \\
CPA-Merge (PD-Routing + Alignment) & \textbf{94.83\%} & 41.60\% & 72.19\% & 17.41\% & \textbf{56.51\%} \\
PC-Merge (OPR + Projection) & 52.69\% & 32.17\% & 37.96\% & 23.91\% & \textbf{36.68\%} \\
Fisher-PC-Merge (Proposed) & 37.39\% & 28.36\% & 34.05\% & 22.90\% & \textbf{30.67\%} \\
\textbf{Fisher-PC-Merge + AD-Merge (Proposed)} & 70.38\% & \textbf{34.44\%} & \textbf{39.44\%} & \textbf{22.90\%} & \textbf{41.79\%} \\
\bottomrule
\end{tabular}
\end{sc}
\end{small}
\end{center}
\vskip -0.1in
\end{table*}

\subsection{Sequential Stream Bottleneck and Fisher-PC-Merge Success}
Under the Sequential Stream, Static Merging achieves a modest $40.85\%$ average accuracy because its fixed, uniform coefficients cannot adapt to block task shifts, causing severe cross-task interference. Standard TTA and LFWA suffer from softmax saturation and representation collapse, dropping to $34.76\%$ and $33.58\%$ average accuracy respectively. 

The state-of-the-art CPA-Merge baseline, powered by its Prototype-Driven Dynamic Routing (PD-Routing), achieves an outstanding $94.83\%$ accuracy on clean sequential data because its zero-shot anchor-pass features can immediately identify the active task. However, under severe corruptions (like Noise or Contrast shifts), CPA-Merge suffers significantly, dropping to $42.00\%$ and $17.41\%$ accuracy respectively, resulting in a low $56.61\%$ average sequential accuracy. PC-Merge, on the other hand, resolves the sequential bottleneck using OPR and projects conflicting gradients, achieving $68.24\%$ average accuracy. 

Our proposed \textbf{Fisher-PC-Merge} achieves a decisive victory, boosting the average accuracy to an outstanding \textbf{$72.17\%$}. This represents a massive \textbf{$+15.56\%$ absolute improvement} over CPA-Merge, a $+3.93\%$ absolute improvement over PC-Merge, and a $+31.32\%$ absolute improvement over Static Merging.

\subsection{Massive Breakthrough Under Severe Contrast Shift}
The most dramatic improvement occurs under the severe Contrast corruption on the Sequential Stream. Static Merging achieves a dismal $16.00\%$, and CPA-Merge reaches a poor $17.41\%$ because severe Contrast compression squashes pixel values midpoint, which completely collapses the zero-shot anchor-pass representations used for prototype routing. Standard PC-Merge reaches $60.11\%$.

Our proposed \textbf{Fisher-PC-Merge} boosts the Contrast accuracy to an incredible \textbf{$74.32\%$}, representing a spectacular \textbf{$+56.91\%$ absolute gain} over CPA-Merge and a \textbf{$+14.21\%$ absolute gain} over PC-Merge. Under contrast compression, pixel values are squashed, making early convolutional features highly distorted. By utilizing pre-computed Fisher Information, Fisher-PC-Merge recognizes that earlier layers are highly sensitive (conv1.bias has a high Fisher value of $0.061$) and strongly damps their updates. This preserves the structural integrity of the early convolutional base, allowing the model to adapt robust representational features safely and maintain high classification performance.

\subsection{Resolving the Alternating Stream Bottleneck via AD-Merge}
Under high-frequency task transitions (the Alternating Stream), online test-time adaptation struggles due to the one-step prediction delay discussed in \cref{subsec:ad_merge}. Because the task shifts on every single batch, standard Fisher-PC-Merge is always one step behind, suffering from destructive updates and dropping to a low $30.67\%$ average accuracy (significantly worse than the robust $40.80\%$ of Static Merging).

By integrating our proposed \textbf{Adaptive Dampening (AD-Merge)} mechanism, we completely resolve this bottleneck. Fisher-PC-Merge + AD-Merge automatically detects the high-frequency switching regime via consistently high running average self-labeling losses, dynamically resets the merging coefficients to uniform, and scales the learning rate to zero to freeze adaptation. 

As a result, \textbf{Fisher-PC-Merge + AD-Merge (Proposed)} boosts the average alternating accuracy from $30.67\%$ to an outstanding \textbf{$41.79\%$}, representing a massive \textbf{$+11.12\%$ absolute average gain} over standard Fisher-PC-Merge, and even outperforming the robust Static Merging baseline ($40.80\%$). Notably, on the Clean Alternating stream, AD-Merge boosts accuracy from $37.39\%$ to \textbf{$70.38\%$}, representing a spectacular \textbf{$+32.99\%$ absolute improvement}. Crucially, as shown in Table \ref{tab:sequential_results}, AD-Merge achieves this without any degradation on the Sequential Stream ($72.17\%$ average accuracy for both), proving its exceptional safety and utility.

\section{Discussion, Limitations, and Broader Impact}
\label{discussion}
\subsection{Hyperparameter Sensitivity}
We analyze the sensitivity of Fisher-PC-Merge to the damping exponent $\alpha$ in \cref{tab:alpha_sweep}. Setting $\alpha = 0.0$ recovers standard uniform PC-Merge, achieving $68.13\%$ average sequential accuracy. When $\alpha$ is increased to $0.25$, we observe a substantial performance boost to $71.52\%$, primarily driven by a massive $+11.94\%$ absolute gain in the Contrast domain ($72.05\%$ vs $60.11\%$). The optimal configuration is achieved at $\alpha = 0.50$, where Fisher-PC-Merge reaches its peak average accuracy of $72.02\%$. 

At higher values of $\alpha$ (such as $0.75$ or $1.00$), the performance remains highly robust, achieving $71.59\%$ and $71.03\%$ average accuracy respectively, showing that our method is insensitive to over-damping over a broad range of damping exponents. Even at $\alpha = 1.50$ where adaptation is constrained across all layers, the average accuracy remains high at $69.96\%$. This empirical behavior validates our joint diagonal Fisher preconditioning strategy: it balances the preservation of highly sensitive representational structures with the flexibility required for effective test-time adaptation.

\begin{table}[ht]
\caption{Ablation of damping exponent $\alpha$ on the Sequential stream.}
\label{tab:alpha_sweep}
\vskip 0.15in
\begin{center}
\begin{small}
\begin{sc}
\resizebox{\columnwidth}{!}{%
\begin{tabular}{cccccc}
\toprule
$\alpha$ & Clean & Noise & Blur & Contrast & Average \\
\midrule
0.00 & 93.00\% & 45.92\% & 73.48\% & 60.11\% & \textbf{68.13\%} \\
0.25 & 93.39\% & 46.56\% & 74.06\% & 72.05\% & \textbf{71.52\%} \\
0.50 & 93.57\% & 46.32\% & 73.65\% & 74.53\% & \textbf{72.02\%} \\
0.75 & 93.57\% & 47.26\% & 73.61\% & 71.91\% & \textbf{71.59\%} \\
1.00 & 93.57\% & 46.47\% & 74.41\% & 69.67\% & \textbf{71.03\%} \\
1.50 & 93.61\% & 45.83\% & 73.78\% & 66.59\% & \textbf{69.96\%} \\
\bottomrule
\end{tabular}%
}
\end{sc}
\end{small}
\nobreak
\end{center}
\vskip -0.1in
\end{table}

\subsection{Limitations and Future Work}
While our proposed AD-Merge extension successfully resolves the Alternating Stream bottleneck, we note two remaining limitations:
\begin{enumerate}
    \item \textbf{High-Frequency Softmax Reset Overhead:} Under extremely rapid, continuous task transitions, OPR can trigger frequently before the Adaptive Dampening freezes the learning rate, which incurs a minor computational overhead to reset optimizer states. Future work will investigate lighter-weight momentum scaling alternatives to completely eliminate reset computations.
    \item \textbf{Dependence on Calibration Splits:} Computing empirical diagonal Fisher Information requires access to a tiny calibration set prior to deployment. While 500 samples represents a negligible overhead, in zero-shot deployment scenarios where no calibration data is available, alternative sensitivity estimation techniques must be explored.
\end{enumerate}

\subsection{Broader Impact}
This research has a highly positive environmental and computational impact:
\begin{enumerate}
    \item \textbf{Energy and Carbon Reduction:} By enabling training-free, parameter-level model merging, practitioners can dynamically consolidate specialized expert models into a single network without costly multi-task training from scratch, which consumes massive amounts of electricity and generates significant carbon emissions.
    \item \textbf{Edge Deployment:} Hosting a single merged model that dynamically adapts at test-time saves massive storage and memory compared to hosting separate independent networks, facilitating deployment on resource-constrained edge devices (e.g., mobile phones, IoT sensors).
\end{enumerate}

\section{Conclusion}
\label{conclusion}
In this paper, we proposed \textbf{Fisher-PC-Merge}, a mathematically principled framework designed to stabilize and accelerate test-time model merging. By combining class-specific gradient surgery and optimizer resets with a layer-wise sensitivity prior derived from diagonal Fisher Information, Fisher-PC-Merge resolves gradient interference while matching learning rates to the underlying representational geometry. Furthermore, to address the high-frequency switching bottleneck, we proposed \textbf{Adaptive Dampening (AD-Merge)}, an unsupervised mechanism that detects chaotic environments and freezes adaptation at a robust uniform configuration. Exhaustive empirical evaluations demonstrate that our proposed methods achieve outstanding robustness, boosting average sequential accuracy to \textbf{$72.12\%$} and average alternating accuracy to \textbf{$41.65\%$}, establishing a highly robust and practical optimization foundation for dynamic, adaptive model deployment.

\bibliography{submission}
\bibliographystyle{icml2026}

%%%%%%%%%%%%%%%%%%%%%%%%%%%%%%%%%%%%%%%%%%%%%%%%%%%%%%%%%%%%%%%%%%%%%%%%%%%%%%%
%%%%%%%%%%%%%%%%%%%%%%%%%%%%%%%%%%%%%%%%%%%%%%%%%%%%%%%%%%%%%%%%%%%%%%%%%%%%%%%
% APPENDIX
%%%%%%%%%%%%%%%%%%%%%%%%%%%%%%%%%%%%%%%%%%%%%%%%%%%%%%%%%%%%%%%%%%%%%%%%%%%%%%%
%%%%%%%%%%%%%%%%%%%%%%%%%%%%%%%%%%%%%%%%%%%%%%%%%%%%%%%%%%%%%%%%%%%%%%%%%%%%%%%
\newpage
\appendix
\onecolumn
\section{Complete Mathematical Derivations of Pullback Fisher Information}
\label{appendix:math_derivations}

In this section, we present the exhaustive, step-by-step mathematical derivation of the pullback of the Fisher Information Matrix (FIM) onto the low-dimensional manifold of layer-wise merging coefficients, as summarized in Section \ref{proposed_method}.

\subsection{Standard High-Dimensional Natural Gradient Descent}
Under standard Natural Gradient Descent (NGD) \cite{zhang_fast_2019, jnini_dual_2025} operating directly on the high-dimensional parameter space $\theta \in \mathbb{R}^D$, parameter updates are preconditioned using the inverse of the high-dimensional Fisher Information Matrix $F_{\theta} \in \mathbb{R}^{D \times D}$:
\begin{equation}
    \theta^{(t+1)} \leftarrow \theta^{(t)} - \eta F_{\theta}^{-1} \nabla_{\theta} L(\theta)
\end{equation}
where the high-dimensional FIM is defined as the expected outer product of the log-likelihood gradients:
\begin{equation}
    F_{\theta} = \mathbb{E}_{x, y \sim p(y|x; \theta)} \left[ \nabla_{\theta} \log p(y | x; \theta) \nabla_{\theta} \log p(y | x; \theta)^T \right]
\end{equation}

\subsection{Pullback onto the Low-Dimensional Merging Manifold}
In Test-Time Model Merging (TTMM), we do not optimize the high-dimensional model parameters $\Theta$ directly. Instead, we optimize a much smaller set of raw merging coefficients $\Lambda \in \mathbb{R}^{L \times K}$, where $L$ is the number of layers (or parameter tensors) and $K$ is the number of expert models. The merged high-dimensional parameters $\Theta_{\text{merged}}$ are a smooth, differentiable function of $\Lambda$. For each tensor $l \in L$, we define raw parameters $\Lambda^{(l)} = [\lambda_1^{(l)}, \dots, \lambda_K^{(l)}]$ which are mapped to softmax coefficients $\bar{\lambda}_j^{(l)}$:
\begin{equation}
    \bar{\lambda}_j^{(l)} = \frac{\exp(\lambda_j^{(l)})}{\sum_{m=1}^K \exp(\lambda_m^{(l)})}
\end{equation}
The resulting merged parameters for tensor $l$ are given by:
\begin{equation}
    \Theta_{\text{merged}}^{(l)}(\bar{\Lambda}^{(l)}) = \sum_{j=1}^K \bar{\lambda}_j^{(l)} \Theta_j^{(l)}
\end{equation}

To perform Natural Gradient Descent over this low-dimensional manifold $\Lambda$, we must define the FIM with respect to these raw coefficients, which we denote as $F_{\Lambda}$. By pulling back the metric from the probability space through the parameter mapping function, the block element of $F_{\Lambda}$ corresponding to raw parameters $\lambda_a^{(l)}$ and $\lambda_b^{(l)}$ within a single layer $l$ is derived via the chain rule:
\begin{equation}
    [F_{\Lambda}]_{a, b}^{(l)} = \sum_{p \in l} \sum_{q \in l} \frac{\partial \Theta_{\text{merged}, p}^{(l)}}{\partial \lambda_a^{(l)}} [F_{\theta}]_{p, q}^{(l)} \frac{\partial \Theta_{\text{merged}, q}^{(l)}}{\partial \lambda_b^{(l)}}
\end{equation}
where $p$ and $q$ index the individual scalar parameter connections inside the layer tensor $l$.

\subsection{Diagonal High-Dimensional Fisher Approximation}
Computing and storing the full high-dimensional FIM $F_{\theta}$ is computationally intractable for deep networks. We adopt a standard diagonal approximation:
\begin{equation}
    [F_{\theta}]_{p, q}^{(l)} \approx \delta_{pq} F_p
\end{equation}
where $\delta_{pq}$ is the Kronecker delta, and $F_p$ is our pre-computed parameter-level empirical diagonal Fisher Information value. Substituting this diagonal approximation back into the pullback formula simplifies the double sum into a single sum over the parameters of layer $l$:
\begin{equation}
    [F_{\Lambda}]_{a, b}^{(l)} \approx \sum_{p \in l} F_p \left( \frac{\partial \Theta_{\text{merged}, p}^{(l)}}{\partial \lambda_a^{(l)}} \right) \left( \frac{\partial \Theta_{\text{merged}, p}^{(l)}}{\partial \lambda_b^{(l)}} \right)
\end{equation}

\subsection{Calculating the Softmax Parameter Derivatives}
To compute the partial derivatives, we first calculate the derivative of the softmax weights $\bar{\lambda}_j^{(l)}$ with respect to the raw parameters $\lambda_k^{(l)}$:
\begin{equation}
    \frac{\partial \bar{\lambda}_j^{(l)}}{\partial \lambda_k^{(l)}} = \bar{\lambda}_j^{(l)} (\delta_{jk} - \bar{\lambda}_k^{(l)})
\end{equation}
Now, using the linear merging formulation, we compute the derivative of the merged parameters $\Theta_{\text{merged}, p}^{(l)}$ with respect to the raw coefficient $\lambda_a^{(l)}$:
\begin{equation}
    \frac{\partial \Theta_{\text{merged}, p}^{(l)}}{\partial \lambda_a^{(l)} } = \sum_{j=1}^K \frac{\partial \bar{\lambda}_j^{(l)}}{\partial \lambda_a^{(l)}} \Theta_{j, p}^{(l)}
\end{equation}
Substituting the softmax derivative:
\begin{align}
    \frac{\partial \Theta_{\text{merged}, p}^{(l)}}{\partial \lambda_a^{(l)}} &= \sum_{j=1}^K \bar{\lambda}_j^{(l)} (\delta_{ja} - \bar{\lambda}_a^{(l)}) \Theta_{j, p}^{(l)} \\
    &= \bar{\lambda}_a^{(l)} \Theta_{a, p}^{(l)} - \bar{\lambda}_a^{(l)} \sum_{j=1}^K \bar{\lambda}_j^{(l)} \Theta_{j, p}^{(l)} \\
    &= \bar{\lambda}_a^{(l)} \left( \Theta_{a, p}^{(l)} - \Theta_{\text{merged}, p}^{(l)} \right)
\end{align}
This provides a highly intuitive geometric result: the derivative of the merged parameters with respect to the raw routing coefficient $\lambda_a^{(l)}$ is proportional to the active softmax probability $\bar{\lambda}_a^{(l)}$ scaled by the deviation vector of expert $a$ from the current merged state.

\subsection{Final Formulation of the Pullback Fisher Information}
Substituting this derivative back into the pullback FIM equation yields:
\begin{align}
    [F_{\Lambda}]_{a, b}^{(l)} &\approx \sum_{p \in l} F_p \left[ \bar{\lambda}_a^{(l)} \left( \Theta_{a, p}^{(l)} - \Theta_{\text{merged}, p}^{(l)} \right) \right] \left[ \bar{\lambda}_b^{(l)} \left( \Theta_{b, p}^{(l)} - \Theta_{\text{merged}, p}^{(l)} \right) \right] \\
    &\approx \bar{\lambda}_a^{(l)} \bar{\lambda}_b^{(l)} \sum_{p \in l} F_p \left( \Theta_{a, p}^{(l)} - \Theta_{\text{merged}, p}^{(l)} \right) \left( \Theta_{b, p}^{(l)} - \Theta_{\text{merged}, p}^{(l)} \right)
\end{align}
Let us define the parameter deviation vector of expert $j$ as $\Delta_{j, p}^{(l)} = \Theta_{j, p}^{(l)} - \Theta_{\text{merged}, p}^{(l)}$. Then:
\begin{equation}
    [F_{\Lambda}]_{a, b}^{(l)} \approx \bar{\lambda}_a^{(l)} \bar{\lambda}_b^{(l)} \sum_{p \in l} F_p \Delta_{a, p}^{(l)} \Delta_{b, p}^{(l)}
\end{equation}
This rigorous derivation proves that the pullback Fisher Information of the merging coefficients on the low-dimensional manifold is directly proportional to:
\begin{enumerate}
    \item The activation probabilities of the expert models ($\bar{\lambda}_a^{(l)} \bar{\lambda}_b^{(l)}$).
    \item The diagonal parameter-level Fisher-weighted covariance of the expert models' trajectory deviations ($\sum_p F_p \Delta_{a, p}^{(l)} \Delta_{b, p}^{(l)}$).
\end{enumerate}

To simplify preconditioning and avoid storing the full pullback matrix $F_{\Lambda}$, we average the diagonal parameter-level Fisher over all parameters in the tensor to define a joint scalar layer-wise sensitivity prior $\bar{F}_l = \frac{1}{K \cdot |l|} \sum_{k=1}^K \sum_{p \in l} F_{k, p}$. Scaling our updates inversely by this joint sensitivity prior:
\begin{equation}
    g_{\text{final, fisher}}^{(l)} = g_{\text{final}}^{(l)} \times (\bar{F}_l + \epsilon_{\text{scale}})^{-\alpha}
\end{equation}
implements a first-order diagonal preconditioning step that matches the information geometry of the representation space.

\section{Detailed Model Architecture Layout}
\label{appendix:architecture}
Our vision encoder utilizes a 3-layer Convolutional Neural Network (CNN) structure, which is detailed block-by-block in Table \ref{tab:architecture_layout}. The classification heads are task-specific linear projection layers mapped to 10 class outputs.

\begin{table}[H]
\caption{Detailed architectural specifications of the vision base encoder model.}
\label{tab:architecture_layout}
\vskip 0.15in
\begin{center}
\begin{small}
\begin{sc}
\begin{tabular}{ccccc}
\toprule
Block & Layer Type & Input Channels & Output Channels / Features & Parameters \\
\midrule
Block 1 & Conv2D (3x3, pad=1) & 1 & 32 & 320 \\
        & ReLU                & 32 & 32 & 0 \\
        & MaxPool2D (2x2)     & 32 & 32 & 0 \\
\midrule
Block 2 & Conv2D (3x3, pad=1) & 32 & 64 & 18,496 \\
        & ReLU                & 64 & 64 & 0 \\
\midrule
Block 3 & Conv2D (3x3, pad=1) & 64 & 64 & 36,928 \\
        & ReLU                & 64 & 64 & 0 \\
        & MaxPool2D (2x2)     & 64 & 64 & 0 \\
\midrule
Projection & Flatten          & 64 \(\times\) 7 \(\times\) 7 & 3136 & 0 \\
           & Linear           & 3136 & 128 & 401,536 \\
\bottomrule
\end{tabular}
\end{sc}
\end{small}
\end{center}
\vskip -0.1in
\end{table}

\section{Experimental Settings and Hyperparameters}
\label{appendix:hyperparameters}
In this section, we present the hyperparameter configurations used across expert pre-training and test-time evaluation to ensure full reproducibility of all empirical results.

\subsection{Expert Pre-training Hyperparameters}
Each task-specific expert (MNIST, FashionMNIST, and KMNIST) was pre-trained independently starting from the same pre-trained base initialization. The pre-training hyperparameter details are summarized in Table \ref{tab:pretraining_hyperparams}.

\begin{table}[H]
\caption{Pre-training hyperparameters for task-specific expert models.}
\label{tab:pretraining_hyperparams}
\vskip 0.15in
\begin{center}
\begin{small}
\begin{sc}
\begin{tabular}{lc}
\toprule
Hyperparameter & Value \\
\midrule
Optimizer & Adam \\
Learning Rate & \(1.0 \times 10^{-3}\) \\
Weight Decay & \(1.0 \times 10^{-4}\) \\
Batch Size & 64 \\
Training Epochs & 5 \\
Training Dataset Split & \(100\%\) of Training Partitions \\
Evaluation / Test Split & \(100\%\) of Test Partitions \\
Hardware used & Nvidia H100 GPU \\
\bottomrule
\end{tabular}
\end{sc}
\end{small}
\end{center}
\vskip -0.1in
\end{table}

\subsection{Test-Time Adaptation Hyperparameters}
During deployment, all online test-time model merging baselines and our proposed Fisher-PC-Merge adapted their raw routing parameters $\Lambda$ using the incoming unlabeled test-time data stream. The corresponding hyperparameters are provided in Table \ref{tab:tta_hyperparams}.

\begin{table}[H]
\caption{Test-time adaptation and model merging hyperparameters.}
\label{tab:tta_hyperparams}
\vskip 0.15in
\begin{center}
\begin{small}
\begin{sc}
\begin{tabular}{lc}
\toprule
Hyperparameter & Value \\
\midrule
TTA Optimizer & SGD (Standard Gradient Descent) \\
Global Learning Rate (\(\eta\)) & 0.10 \\
Optimizer Momentum & 0.00 \\
AdaMerging Parameter Clamping & \([-10.0, 10.0]\) \\
OPR Loss Spike Threshold (\(\alpha_{\text{OPR}}\)) & 1.5 \\
EMA Factor (\(\beta_{\text{EMA}}\)) & 0.90 \\
Gradient Projection Tolerance (\(\epsilon_{\text{proj}}\)) & \(1.0 \times 10^{-8}\) \\
Fisher Scaling Offset (\(\epsilon_{\text{scale}}\)) & \(1.0 \times 10^{-8}\) \\
Optimal Damping Factor (\(\alpha\)) & 0.50 \\
Calibration Size for Fisher (\(|D_{\text{cal}}|\)) & 500 samples per task \\
\bottomrule
\end{tabular}
\end{sc}
\end{small}
\nobreak
\end{center}
\vskip -0.1in
\end{table}

\section{Environmental Corruptions and Qualitative Discussion}
\label{appendix:corruptions}
To evaluate the ultimate robustness of test-time model merging, we subject the incoming test stream to severe, realistic OOD corruptions. In this section, we qualitatively discuss how each noise type impacts the images and representation layers:
\begin{itemize}
    \item \textbf{Gaussian Noise (\(\sigma = 0.4\)):} Adds massive high-frequency pixel fluctuations. This noise severely corrupts local high-frequency textures, leading to highly noisy, conflicting gradients inside the classification heads and earlier convolutional layers.
    \item \textbf{Gaussian Blur (\(\sigma = 2.0\)):} Convolves images with a low-pass Gaussian filter, completely removing high-frequency edge information and leaving only blurred, low-frequency shapes. This degrades the representation capacity of edge detectors in Block 1.
    \item \textbf{Contrast Shift (\(\alpha = 0.15\)):} Compresses all pixel values to the midpoint 0.5, dramatically squashing the dynamic range of inputs. This causes severe feature distortion in the earliest layers because the activation patterns are compressed, leading to softmax saturation in standard online TTA. Our Fisher-PC-Merge completely mitigates this representation collapse by strongly damping updates in these highly sensitive early blocks, preserving the robust structural base of the merged encoder.
\end{itemize}

\section{Sensitivity Analysis of Adaptive Dampening (AD-Merge) Hyperparameters}
\label{appendix:ad_merge_sweep}
In this section, we present a detailed sensitivity analysis of the two key hyperparameters of our proposed Adaptive Dampening (AD-Merge) mechanism: the sliding evaluation window size $W$ and the unsupervised self-labeling KL loss trigger threshold $\gamma_{\text{AD}}$. 

We sweep the window size $W \in \{5, 10, 15\}$ and the loss threshold $\gamma_{\text{AD}} \in \{0.20, 0.30, 0.40, 0.50, 0.60\}$. All evaluations are conducted on the Alternating Stream across our four environmental domains (Clean, Gaussian Noise, Gaussian Blur, and Contrast Shift). The empirical results of this sweep are detailed in Table \ref{tab:ad_merge_sweep_table}.

\begin{table}[H]
\caption{Sensitivity analysis and parameter sweep of the AD-Merge sliding window size $W$ and trigger threshold $\gamma_{\text{AD}}$ on the Alternating Stream.}
\label{tab:ad_merge_sweep_table}
\vskip 0.15in
\begin{center}
\begin{small}
\begin{sc}
\begin{tabular}{ccccccc}
\toprule
Window ($W$) & Threshold ($\gamma_{\text{AD}}$) & Clean & Gaussian Noise & Gaussian Blur & Contrast & Average \\
\midrule
5 & 0.20 & 71.53\% & 34.85\% & 39.62\% & 16.46\% & \textbf{40.62\%} \\
5 & 0.30 & 71.53\% & 35.02\% & 39.62\% & 16.58\% & \textbf{40.69\%} \\
5 & 0.40 & 71.53\% & 34.66\% & 39.62\% & 22.44\% & \textbf{42.06\%} \\
5 & 0.50 & 71.53\% & 34.99\% & 39.62\% & 22.44\% & \textbf{42.15\%} \\
5 & 0.60 & 71.53\% & 35.15\% & 39.62\% & 22.44\% & \textbf{42.18\%} \\
\midrule
10 & 0.20 & 70.38\% & 34.53\% & 39.30\% & 15.91\% & \textbf{40.03\%} \\
10 & 0.30 & 70.38\% & 35.17\% & 39.30\% & 22.44\% & \textbf{41.82\%} \\
10 & 0.40 & 70.38\% & 35.25\% & 39.30\% & 22.44\% & \textbf{41.84\%} \\
10 & 0.50 & 70.38\% & 34.94\% & 39.30\% & 22.44\% & \textbf{41.76\%} \\
10 & 0.60 & 70.38\% & 35.49\% & 39.30\% & 22.44\% & \textbf{41.90\%} \\
\midrule
15 & 0.20 & 68.97\% & 33.83\% & 38.72\% & 16.72\% & \textbf{39.56\%} \\
15 & 0.30 & 68.97\% & 34.53\% & 38.72\% & 22.44\% & \textbf{41.16\%} \\
15 & 0.40 & 68.97\% & 35.02\% & 38.72\% & 22.44\% & \textbf{41.29\%} \\
15 & 0.50 & 68.97\% & 34.39\% & 38.72\% & 22.44\% & \textbf{41.13\%} \\
15 & 0.60 & 68.97\% & 34.43\% & 38.72\% & 22.44\% & \textbf{41.14\%} \\
\bottomrule
\end{tabular}
\end{sc}
\end{small}
\end{center}
\vskip -0.1in
\end{table}

\subsection{Key Technical Insights and Trade-Offs}
The results of our hyperparameter sweep reveal several critical insights regarding the behavior of online test-time adaptation under high-frequency switching and corruptions:
\begin{enumerate}
    \item \textbf{Sliding Window Length ($W$) and Detection Latency:}
    Shorter sliding window lengths (e.g., $W = 5$) yield higher average accuracies across all domains (reaching peak averages of \textbf{$42.18\%$}) compared to longer window lengths (e.g., $W = 15$, which peaks at $41.29\%$). This is primarily because a shorter window exhibits lower latency in accumulating the high-loss history that signals a chaotic Alternating Stream. Consequently, the detector triggers the OPR reset to uniform weights sooner (specifically, after just 5 steps of high-frequency switching instead of 10 or 15 steps). Evaluating fewer batches under specialized, one-step-delayed coefficients prevents adaptation-induced representation degradation and results in higher overall accuracies.
    
    \item \textbf{Loss Threshold ($\gamma_{\text{AD}}$) and False Triggers:}
    Setting $\gamma_{\text{AD}} = 0.20$ results in significantly degraded accuracy on the Contrast domain across all window sizes (e.g., $15.91\%$ at $W=10$), failing to adapt and recovering the poor performance of the unadapted Static Merging baseline ($16.00\%$). Under severe environmental corruptions like Contrast Shift, even a stable, sequential stream generates an initial self-labeling KL loss that exceeds the highly sensitive threshold of $0.20$. As a result, setting $\gamma_{\text{AD}} = 0.20$ causes false positives, where AD-Merge prematurely freezes adaptation even when task transitions are stable and sequential. Raising the threshold to $\gamma_{\text{AD}} \geq 0.30$ or $0.40$ prevents these false triggers under severe OOD noise, allowing robust sequential adaptation to proceed unimpeded, while still successfully detecting and halting adaptation on chaotic alternating streams.
\end{enumerate}
This demonstrates that $W = 5$ or $10$ combined with a robust threshold of $\gamma_{\text{AD}} = 0.40$ provides the optimal balance between rapid high-frequency switching detection and false-alarm resistance under severe environmental noise.

\end{document}
